\clearpage
\begingroup
\let\Z\relax
\newcommand{\Z}{\mathbb Z}
\let\N\relax
\newcommand{\N}{\mathbb N}
\let\rad\relax
\newcommand{\rad}{\operatorname{rad}}
\let\eps\relax
\newcommand{\eps}{\varepsilon}
\section*{Integrated checkpoint: First-appearance depth and the norm-seven orbit}
\phantomsection\label{checkpoint:rank}
\addcontentsline{toc}{section}{Integrated checkpoint: First-appearance depth and the norm-seven orbit}
\noindent\textit{Edited proof consolidation; provenance and historical verification boundaries are stated below.}\medskip

\section{Provenance and the fixed arithmetic object}
This is a proof consolidation of the preceding rank/depth supplement, not a byte-identical reprint of its historical runtime narrative. The original conversation source is identified by Git blob \texttt{69dc72153b9dcdc58e7b512b20a57cb46e03a733}. The local laws, truncation arguments, modular certificate and fixed-angle application are restated explicitly below. The complex logarithmic-form input is cited in a directly checked form. Neither the original nor this edition proves standard abc.

Let $\zeta=(1+\sqrt{-3})/2$, let $z_n=(2+\zeta)^n=(x_n,y_n)$, and define
\[
 N(x,y)=x^2+xy+y^2,\quad P(x,y)=xy(x+y),\quad H(x,y)=\max(|x|,|y|,|x+y|).
\]
Then
\[
 z_{n+1}=(2x_n-y_n,x_n+3y_n),\qquad N(z_n)=7^n.
\]
These pairs are primitive: a common prime after one step must divide both $7x_n$ and $7y_n$ and hence can only be seven; but for $n\ge1$, $z_n\equiv5^{n-1}(2,1)\pmod7$, excluding it. Unit rotation supplies a positive triple for every $n\ge1$. A primitive zero-boundary pair is a unit, incompatible with norm $7^n>1$.

Put $H_n=H(z_n)$, $\mathscr R_n=\rad(|P(z_n)|)$. Cubic expansion yields
\[
 P(z_n)=6U_n,\qquad U_0=0,\quad U_1=1,\quad U_{n+2}=20U_{n+1}-343U_n.
\]
Equivalently, with $\alpha=10+9\sqrt{-3}$ and $\beta=10-9\sqrt{-3}$,
\[
 U_n=\frac{\alpha^n-\beta^n}{\alpha-\beta}.
\]
The earlier full boundary gcd theorem gives $\gcd(|U_m|,|U_n|)=|U_{\gcd(m,n)}|$. In particular a prime with least positive appearance index $d_p$ divides $U_n$ exactly when $d_p\mid n$.

\section{The rank and valuation laws, including exceptional primes}
For every prime that appears, define $s_p=v_p(U_{d_p})$.
\begin{theorem}[Complete local laws]\label{rank:thm:local}
Seven never divides $U_n$ for $n\ge1$. Every other prime has a first index, and
\[
 v_p(U_n)=0\quad(d_p\nmid n),\qquad
 v_p(U_n)=s_p+v_p(n/d_p)\quad(d_p\mid n).
\]
At the exceptional primes,
\[
 d_2=2,\ s_2=2,\quad
 v_2(U_n)=\begin{cases}0&n\text{ odd},\\v_2(n)+1&n\text{ even},\end{cases}
 \qquad d_3=3,\ s_3=1,\quad v_3(U_n)=v_3(n).
\]
For $p\ge5$, $p\ne7$, set $\chi_p=1$ when $p\equiv1\pmod3$ and $\chi_p=-1$ when $p\equiv2\pmod3$. Then
\[
 d_p\mid(p-\chi_p)/3,\qquad p\ge3d_p-1.
\]
\end{theorem}
\begin{proof}
Modulo seven, the recurrence gives $U_n\equiv20^{n-1}\ne0$. For $p\ne2,3,7$ work in the unramified quadratic algebra with $s^2=-3$, which is either a field or a product of two fields. Let $\gamma=(5+s)/2$, of norm seven. Since $\gamma^3=\alpha$, the quotient $\alpha/\beta=(\gamma/\bar\gamma)^3$ is a cube in the norm-one group. Frobenius is identity or conjugation according as $\chi_p$ is one or minus one. Thus $(\gamma/\bar\gamma)^{p-\chi_p}=1$. As three divides $p-\chi_p$, the order of $\alpha/\beta$ divides $(p-\chi_p)/3$. The denominator $\alpha-\beta$ and $\beta$ are units at $p$, so this order is precisely the first index of $U_n$.
In each unramified component, write $(\alpha/\beta)^{d_p}=1+p^{s_p}u$ with $u$ a unit. For a positive multiplier coprime to $p$, binomial expansion leaves the first term at valuation $s_p$. Raising to the $p$-th power increases it by exactly one, since the first binomial term has uniquely least valuation. Iteration proves $s_p+v_p(n/d_p)$. The component valuations agree with the valuation of the rational integer $U_n$, since the Binet denominator and $\beta$ are units.

At three use the ramified field $\Q_3(\sqrt{-3})$, with valuation extending $v_3$ on $\Q$. Here $v_3((\alpha/\beta)-1)=5/2$. The same binomial argument is valid at this depth: the first term remains uniquely of least valuation, giving $v_3((\alpha/\beta)^n-1)=5/2+v_3(n)$. Subtracting the denominator valuation gives $v_3(U_n)=v_3(n)$.
At two, the recurrence shows $U_n$ is odd exactly when $n$ is odd. Put $V_n=\alpha^n+\beta^n$. Then $U_{2n}=U_nV_n$, and $V_0=2,V_1=20$ satisfy the same recurrence. Modulo eight it has period two, with $V_n\equiv4$ for odd $n$ and $2$ for even $n$. Thus $v_2(V_n)=2$ in the first case and one in the second. Starting at odd $n$ and doubling repeatedly proves the displayed two-adic law. The first indices and exponents follow, including $U_2=20$ and $U_3=57$.
\end{proof}

\section{Exact bookkeeping and absorption through depth two}
Define
\[
 r_n=\rad(|U_n|),\quad
 \mathcal L_n=\prod_{d_p\mid n}p^{v_p(n/d_p)},\quad
 \mathcal D_n=\prod_{d_p\mid n}p^{(s_p-2)_+},\quad
 F_n=\prod_{d_p\mid n}p^{(s_p-3)_+}.
\]
Every product is over the actual finite prime support of $U_n$; $u_+=\max(u,0)$. Theorem~\ref{rank:thm:local} implies $\mathcal L_n\mid n$ and
\[
 |U_n|=\mathcal L_n\prod_{d_p\mid n}p^{s_p}.
\]
In particular
\[
 |U_n|\le n r_n^2\mathcal D_n,\qquad
 |U_n|\le n r_n^3 F_n.
\]
The index-lifting factor is not discarded: it is bounded by $n$ because $v_p(n/d_p)\le v_p(n)$ for every supported prime.

For positive rotated coordinates $x,y$, let $H=x+y$. Since $xy\ge x+y-1$ and $xy\ge1$, $H\le2xy$, and therefore $H^2\le2xyH=2P$. Also
\[
 P+1-N=(x+y+1)(x-1)(y-1)\ge0.
\]
Consequently the following is unconditional:
\begin{theorem}[Elementary depth-two absorption]\label{rank:thm:two}
For every $n\ge1$,
\[
 H_n^2\le12n\mathscr R_n^2\mathcal D_n.
\]
\end{theorem}
\begin{proof}
Use $H_n^2\le2|P(z_n)|=12|U_n|$ and $r_n\le\mathscr R_n$ in the preceding inequality.
\end{proof}
Thus first squares do not occur in the remaining excess $\mathcal D_n$, and ordinary lifting costs only a factor $n$. This is not a proof that the remaining product is small.

\section{Exact cyclotomic blocks}
For $m\ge1$ define the positive rational product
\[
 A_m=\prod_{d\mid m}|U_d|^{\mu(m/d)}.
\]
All $U_d$ for $d\ge1$ are nonzero. For a prime $p$ with first index $d_p$, its valuation in this product is zero when $d_p\nmid m$. If $m=d_p k$, it is
\[
 \sum_{e\mid k}\mu(k/e)(s_p+v_p(e))
 =\begin{cases}s_p&k=1,\\1&k=p^j,\ j\ge1,\\0&\text{otherwise}.\end{cases}
\]
Indeed the constant term has divisor sum zero except at $k=1$; write $v_p(e)=\sum_{j\ge1}{\bf1}_{p^j\mid e}$ and apply the same divisor cancellation. All valuations are nonnegative, so $A_m$ is an integer. This proves both its first-appearance factorization and the fact that recycled exceptional factors have only exponent one. In particular repeated-prime excess in $A_m$ comes only from first-appearance primes. A bound on how many such primes appear does not by itself bound the first exponents $s_p$.

\section{Complex logarithms absorb first cubes on the fixed orbit}
We use the classical complex logarithmic-form estimate as recorded in Bugeaud, Theorem 1.1, equation (1.2) \cite{rank:bugeaud}: a nonzero integer linear combination of fixed logarithms of fixed algebraic numbers has absolute value bounded below by a negative power of the maximum coefficient. It is an established external theorem, not a Lean axiom and not a real-positive-logarithm theorem applied outside its hypotheses.

\begin{lemma}[Fixed-angle lower bound]\label{rank:lem:angle}
For $\lambda=\alpha/\beta=(-143+180\sqrt{-3})/343$, there are effective fixed constants $c_0>0$, $\kappa\ge0$ such that
\[
 |\lambda^n-1|\ge c_0n^{-\kappa}\qquad(n\ge1).
\]
\end{lemma}
\begin{proof}
The number $\lambda$ has absolute value one, is nonrational, and has quadratic trace $-286/343$, not an integer. Hence it is not an algebraic integer and not a root of unity. Fix $\log\lambda=i\theta$ with $|\theta|\le\pi$, and choose $j\in\Z$ so that $|n\theta-2\pi j|\le\pi$. Then $2|j|\le n+1$, and the form
\[
 \Lambda_n=n\log\lambda-2j\log(-1)
\]
is nonzero. Applying the cited theorem to these two fixed logarithms gives $|\Lambda_n|\ge c_1n^{-\kappa}$ after absorbing finitely many small coefficients. For $|t|\le\pi$, $|e^{it}-1|\ge2|t|/\pi$. This proves the assertion.
\end{proof}

\begin{theorem}[Depth-three absorption]\label{rank:thm:three}
There is a fixed effective $C_0>0$ such that for all $n\ge1$,
\[
 H_n^3\le C_0 n^{\kappa+1}\mathscr R_n^3 F_n.
\]
\end{theorem}
\begin{proof}
Since $|\alpha|=|\beta|=7^{3/2}$ and $|\alpha-\beta|=18\sqrt3$, the Binet formula and Lemma~\ref{rank:lem:angle} give
\[
 |U_n|\ge\frac{c_0}{18\sqrt3}7^{3n/2}n^{-\kappa}.
\]
Norm-height comparison gives $H_n\le(2/\sqrt3)7^{n/2}$. Consequently $H_n^3\le(48/c_0)n^\kappa|U_n|$. Now use $|U_n|\le n r_n^3F_n$ and $r_n\le\mathscr R_n$.
\end{proof}

\begin{corollary}[Uniform restricted class]\label{rank:cor:class}
For every fixed $A\ge0$ and $\eps>0$ there is $K_{A,\eps}>0$ such that
\[
 H_n\le K_{A,\eps}\mathscr R_n^{1+\eps}
\]
whenever $F_n\le n^A$.
\end{corollary}
\begin{proof}
The theorem gives $H_n\le C_0^{1/3}\mathscr R_n n^B$, where $B=(\kappa+1+A)/3$. Since $H_n\ge7^{n/2}$, we have $n\le2\log H_n/\log7$. For $\delta=\eps/(1+\eps)$,
\[
 (\log t)^B\le(B/(e\delta))^B t^\delta\quad(t\ge1).
\]
Substituting, rearranging, and raising to $1/(1-\delta)$ proves the result, with constants depending only on the fixed displayed parameters.
\end{proof}
No proof that this class contains all indices, or any specified infinite family, is asserted. In particular $\log F_n=o(n)$ remains an unproved sufficient estimate for the entire fixed orbit. If $H_n>K\mathscr R_n^{1+\eps}$, the theorem instead forces
\[
 F_n>\frac{K^{3/(1+\eps)}}{C_0n^{\kappa+1}}
             H_n^{3\eps/(1+\eps)}.
\]
Thus any fixed-exponent abc counterfamily within this orbit requires exponentially large first-depth-four excess, after the index contribution is removed. This is a necessary condition, not a construction or an exclusion of such a family.

\section{A genuine first square, checked by exact modular arithmetic}
The prime $p=38629$ has $d_p=12876=2^2\cdot3\cdot29\cdot37$ and $s_p=2$. Trial division by integers through 196 proves primality. In the pair algebra modulo $p$, the element $\alpha=1+18\zeta$ to the full exponent 12876 is $(1,0)$. At exponents $12876/q$ for $q=2,3,29,37$, the results are respectively
\[
 (25926,25406),\quad(38628,1),\quad
 (10564,19690),\quad(23253,26998).
\]
None is scalar, so none has $U$-value zero modulo $p$. The rank divides the full index by Theorem~\ref{rank:thm:local}; any proper divisor divides one of the four tested exponents, so its exact value is 12876.
At modulus $p^2=1492199641$ the full result is $(89232991,0)$; at modulus $p^3=57642179932189$ it is
\[
 (43039603478354,40547540844893),\qquad
 40547540844893=27173p^2,\quad0<27173<p.
\]
The second coordinate is $18U_{12876}$, proving exact first valuation two. These data refute the first-squarefree shortcut, but are absorbed by both of the depth estimates above. They are not an abc counterexample. The current exact replay verifies the second coordinates independently by the linear recurrence, not only by pair repeated squaring.

\section{Formal and global boundary}
The existing \texttt{RankDepth.lean} verifies its explicit finite-list accounting, first-index implication under a stated strong-divisibility hypothesis, elementary boundary bounds and modular certificates. Its successful source run is recorded separately from this ordinary proof edition. The full finite-field order law, ramified and unramified valuation arguments, Binet interpretation, complex logarithmic-form theorem and real-exponent absorption are not all formalized by that module.

The new moving-support supplement studies how to replace the fixed constants in Lemma~\ref{rank:lem:angle} when several norm generators vary. This does not remove the unresolved first-depth product on the fixed orbit, nor the prime-norm terminal and general mixed-factor support problems. No global abc theorem follows from the displayed sufficient classes alone.


\endgroup
